\exercise{Simple  timescale renormalization}{1}
Change the timescale to make the parameter $k$ disappear form the system
\eqn{\dot{x}=k(x-x^2)}


\solution 
We define 
\eq{
\tau = a t.
}
Hence
\eq{
\frac{\partial t}{\partial \tau} = \frac{1}{a}.
}
Therefore
\eq{
\frac{{\rm d}x}{{\rm d}\tau} = \frac{\partial t}{\partial \tau} \frac{{\rm d}x}{{\rm d}t} =\frac{1}{a} k(x-x^2)
}
Now $\tau$ is our new time, and we define
\eq{
\dot{x}=\frac{{\rm d}x}{{\rm d}\tau}  
}
which allows us to write
\eq{
\dot{x}=\frac{k}{a} (x-x^2).
}
Choosing $a=k$, yields the desired equation
\eq{
\dot{x}= (x-x^2)
}
Note we can always set one of the rates in the system to one by means of a timescale renormalization. This just means we measure all other rates realtive to the rate we set to one. So for example if you have to processes which have rate constants $a= 10/{\rm s}$ and $b= 20/{\rm s}$, then you can renormalize time such that $a=1$ and $b=2$ (without the unit seconds). This makes the $a$ disappear and the rate $b$ is now measured in multiples of $a$ instead of seconds.